\documentclass[9pt]{extarticle}
\usepackage[T1]{fontenc}
\usepackage{lmodern,amsmath,amssymb,amsthm,microtype}
\usepackage[margin=0.70in]{geometry}
\usepackage[colorlinks=true,urlcolor=blue,citecolor=blue,linkcolor=blue]{hyperref}
\newtheorem{theorem}{Theorem}
\newtheorem{lemma}{Lemma}
\newtheorem{corollary}{Corollary}
\newcommand{\ord}{\operatorname{ord}}
\newcommand{\len}{\mathcal{L}}
\newcommand{\rev}{\mathcal{R}}
\title{Prime Divisibility of Descending Concatenations in Arbitrary Bases\\and OEIS A077183--A077185}
\author{Qichao Wang\\{\small Hebei University of Technology, Tianjin, China}\\[0.4em]Daoyu Dong\\{\small University of Electronic Science and Technology of China, Chengdu, China}}
\date{October 17, 2026}\linespread{0.97}\setlength{\abovedisplayskip}{6pt}\setlength{\belowdisplayskip}{6pt}\setlength{\abovedisplayshortskip}{3pt}\setlength{\belowdisplayshortskip}{3pt}
\begin{document}
\maketitle
\begin{abstract}
Let $\rev_b(k)$ denote the base-$b$ integer obtained by concatenating
$k,k-1,\ldots,2,1$ in descending order. We prove that a prime $p$
divides at least one such integer if and only if $p\nmid b$, and
that for every prime $p\nmid b$ infinitely many endpoints $k$
exist. For $p\nmid b(b-1)$ the proof gives an explicit construction
with $k< b^{2r+1}$, where $r=\ord_p(b)$, using a nonzero affine
translation over $\mathbb F_p$. In decimal this proves the
2005 positivity conjecture for OEIS A077183 and the associated
A077185 quotient sequence: the only zero terms occur at prime indices
$1$ and $3$. All assertions are independent of finite searches.
\end{abstract}

\section{Definitions and main theorem}
For an integer base $b\ge2$, write $\ell_b(t)$ for the number of
base-$b$ digits of $t\ge1$, and define
\[
\len_b(k)=\sum_{t=1}^{k}\ell_b(t),\qquad
\rev_b(k)=(k)_b\Vert(k-1)_b\Vert\cdots\Vert(1)_b,
\]
where the concatenated digits are interpreted as a single integer.
Set $\rev_b(0)=0$ and $\len_b(0)=0$. Thus
\begin{equation}\label{eq:prepend}
\rev_b(k)=\rev_b(k-1)+k\,b^{\len_b(k-1)} \quad(k\ge1).
\end{equation}
OEIS A077183 asks for the least $k$ such that the $n$-th prime divides
$\rev_{10}(k)$, and records zero when there is no such $k$. Its
companion A077185 is the quotient of this smallest concatenation by
the same prime \cite{oeis183,oeis185,oeis422}. The 2005 conjecture
asserts that A077183 has zero terms only at the indices of $2$ and $5$.

\begin{theorem}[Prime dichotomy in every base]\label{thm:base}
Let $b\ge2$ be an integer and $p$ be a prime. Then
\[
\bigl(\exists k\ge1:\ p\mid\rev_b(k)\bigr)
\quad\Longleftrightarrow\quad p\nmid b.
\]
When $p\nmid b$, infinitely many distinct such $k$ exist. If
$p\nmid b(b-1)$ and $r=\ord_p(b)$, a solution may be chosen satisfying
\begin{equation}\label{eq:bound}
1\le k\le b^{2r}-1+(p-1)r < b^{2r+1}.
\end{equation}
\end{theorem}

\section{A constant, nonzero translation}
\begin{lemma}[Fixed-width period step]\label{lem:shift}
Suppose $p\nmid b(b-1)$ and put $r=\ord_p(b)$.
Let $d\equiv1\pmod r$, $L=b^{d-1}$, and
$A=b^{\len_b(L-1)}\pmod p$. Whenever the integers
$L,L+1,\ldots,L+jr-1$ all have exactly $d$ digits,
\begin{equation}\label{eq:affine}
\rev_b(L-1+jr)
\equiv \rev_b(L-1)+j A\,r(b-1)^{-1}\pmod p.
\end{equation}
The coefficient of $j$ is nonzero modulo $p$.
\end{lemma}
\begin{proof}
Write $k=L+i$, where $0\le i$ and $k$ has $d$ digits.
The exponent in \eqref{eq:prepend} is
\[
\len_b(k-1)=\len_b(L-1)+di.
\]
Because $d\equiv1\pmod r$ and $b^r\equiv1\pmod p$,
we have $b^{di}\equiv b^i\pmod p$. Consequently, after adding
$j r$ new blocks,
\begin{equation}\label{eq:sum}
\rev_b(L-1+jr)-\rev_b(L-1)
\equiv A\sum_{i=0}^{jr-1}(L+i)b^i\pmod p.
\end{equation}
The standard geometric identities give
\begin{equation}\label{eq:geometric}
\sum_{i=0}^{r-1}b^i\equiv0,\qquad
\sum_{i=0}^{r-1}i b^i\equiv\frac{r}{b-1}\pmod p.
\end{equation}
For completeness, $b^r\equiv1$ implies
$\sum_{i=0}^{r-1}b^i=(b^r-1)/(b-1)\equiv0$;
applying a finite geometric derivative gives
\[
\sum_{i=0}^{r-1}i b^i
=\frac{b-rb^r+(r-1)b^{r+1}}{(1-b)^2}
\equiv\frac{r}{b-1}\pmod p.
\]
In each consecutive $r$-term group of \eqref{eq:sum},
$b^{i+r}\equiv b^i$. The starting offset increases by $r$,
but its contribution vanishes because $\sum b^i\equiv0$.
Thus every group contributes exactly $Ar/(b-1)$, proving
\eqref{eq:affine}. Since $p\nmid b$, $A\ne0$; also
$r\mid p-1$ and hence $p\nmid r$. Finally $p\nmid b-1$,
so $Ar/(b-1)\ne0\pmod p$.
\end{proof}

\section{Proof and explicit witness}
\begin{proof}[Proof of Theorem~\ref{thm:base}]
First, if $p\mid b$, then \eqref{eq:prepend} yields
$\rev_b(k)\equiv\rev_b(k-1)\equiv\cdots\equiv1\pmod p$;
there is no solution. If $p\mid b-1$, then $b\equiv1\pmod p$,
so every concatenation obeys
\[
\rev_b(k)\equiv1+2+\cdots+k=\frac{k(k+1)}2\pmod p.
\]
In particular, for every $s\ge1$, the choice $k=2ps$ is a
solution, proving infinitude in this case (also when $p=2$).

It remains to suppose $p\nmid b(b-1)$.
Write $r=\ord_p(b)$, take $d=2r+1$, and set
\[
L=b^{2r},\quad m_0=L-1,\quad
A=b^{\len_b(m_0)}\pmod p,\quad
\Delta=A\,r(b-1)^{-1}\pmod p.
\]
Here $r\mid p-1$, so $\Delta$ is invertible modulo $p$.
Choose $j$ as the least nonnegative residue satisfying
\begin{equation}\label{eq:j}
\boxed{\quad j\equiv-\rev_b(m_0)\Delta^{-1}\pmod p,\quad
0\le j<p.\quad}
\end{equation}
Because $p\mid b^r-1$, we have $p<b^r$; also
$r<(b-1)b^r$ for all $b\ge2,r\ge1$. Therefore
\[
(p-1)r<(b-1)b^{2r},
\]
so the integers $L,\ldots,L-1+(p-1)r$ remain inside the
$d$-digit range $[L,bL-1]$. Applying Lemma~\ref{lem:shift},
\[
\rev_b(m_0+jr)\equiv\rev_b(m_0)+j\Delta\equiv0\pmod p.
\]
The displayed bound \eqref{eq:bound} follows from $j\le p-1$.

For infinitely many further solutions, repeat the construction with
$d=1+tr$ for any sufficiently large integer $t$. The same
constant translation applies, and the $d$-digit window eventually
contains all $p$ candidates. The resulting endpoints have
unbounded digit lengths, hence are distinct.
\end{proof}

\begin{corollary}[Resolution of A077183 and A077185]\label{cor:oeis}
For all $n\ge1$,
\[
\boxed{\mathrm{A077183}(n)>0\iff n\notin\{1,3\}.}
\]
The displayed zero values of A077185 are precisely at indices
$1$ and $3$; all other terms are well-defined positive integers.
For every prime $p\ge7$ the least endpoint in A077183 is bounded
by \eqref{eq:bound} with $b=10$ and $r=\ord_p(10)$.
\end{corollary}
\begin{proof}
Apply Theorem~\ref{thm:base} with $b=10$. The prime divisors
of the base are exactly $2$ and $5$. The exception $p=3$ is
covered by $3\mid10-1$; indeed $\rev_{10}(2)=21$.
The index of prime $2$ is $1$ and that of prime $5$ is $3$.
Since A077185 is the corresponding quotient by a positive prime,
its terms are positive exactly at the nonexceptional indices.
\end{proof}

\section{Verification and perspective}
The accompanying Python script independently reproduces the first
63 published A077183 endpoints and the first 12 A077185 quotient
values; it constructs and verifies the modular witness for primes up
to 101 (and selected additional bases). It never builds the
astronomically large concatenations for these constructive witnesses:
modular geometric sums evaluate the blocks exactly. The symbolic
proof, not the finite calculations, establishes the result.

The proof concerns existence, not a closed form for the \emph{least}
endpoint or the associated quotient. Its mechanism differs from the
ascending or cyclic-digit-rotation Champernowne constructions: the
natural recurrence here prepends an integer block, and the nonzero
period increment comes directly from a weighted geometric sum.

\begin{thebibliography}{9}
\bibitem{oeis183} OEIS Foundation Inc., \emph{A077183: Smallest $k$
for which the reverse concatenation $k,\ldots,1$ is divisible by
prime($n$)}, \url{https://oeis.org/A077183} (accessed October 17, 2026).
\bibitem{oeis185} OEIS Foundation Inc., \emph{A077185},
\url{https://oeis.org/A077185} (accessed October 17, 2026).
\bibitem{oeis422} OEIS Foundation Inc., \emph{A000422: Reverse
concatenation of positive integers},
\url{https://oeis.org/A000422}.
\bibitem{stephan} R. Stephan, \emph{Factors and Primes in Two
Smarandache Sequences}, Smarandache Notions Journal \textbf{9}
(1998), 4--10 (linked from OEIS A077183).
\end{thebibliography}
\end{document}
